\documentclass[11pt]{article}
\usepackage[margin=1in]{geometry}
\usepackage{amsmath, amssymb, amsthm}
\usepackage{mathtools}
\usepackage{hyperref}

\newtheorem{theorem}{Theorem}
\newtheorem{lemma}[theorem]{Lemma}
\newtheorem{corollary}[theorem]{Corollary}
\theoremstyle{definition}
\newtheorem{definition}[theorem]{Definition}
\newtheorem{example}[theorem]{Example}

\title{Sample Spaces}
\author{math-foundations / 17-sample-spaces}
\date{}

\begin{document}
\maketitle

\section*{1. Introduction}

Probability theory begins with a single primitive: a set whose elements are
the possible outcomes of an experiment. This set is the \emph{sample space}.
Every subsequent concept --- events, random variables, expectations,
conditional distributions --- is built on top of this primitive. Before
we can assign probabilities, we must first agree on what we are assigning
probabilities \emph{to}.

\section*{2. Definitions}

\begin{definition}[Sample Space]
A \emph{sample space} is a non-empty set $\Omega$ whose elements represent
the mutually exclusive and collectively exhaustive outcomes of an
experiment or random phenomenon.
\end{definition}

\begin{definition}[Outcome]
An \emph{outcome} is an element $\omega \in \Omega$. It represents one
complete atomic realization of the experiment.
\end{definition}

\begin{definition}[Event]
An \emph{event} is a subset $A \subseteq \Omega$. We say the event $A$
\emph{occurs} on outcome $\omega$ if $\omega \in A$. In the finite case,
the family of all events is the power set
\[
2^{\Omega} \;=\; \{A : A \subseteq \Omega\}.
\]
\end{definition}

\begin{example}[Coin flip]
$\Omega = \{H, T\}$. The events are $\emptyset, \{H\}, \{T\}, \{H, T\}$.
\end{example}

\begin{example}[Die roll]
$\Omega = \{1, 2, 3, 4, 5, 6\}$. The event ``even'' is $\{2, 4, 6\}$.
\end{example}

\begin{example}[Two dice]
$\Omega = \{1, \ldots, 6\}^2$, with $|\Omega| = 36$. The event ``sum is
$7$'' is $\{(1,6), (2,5), (3,4), (4,3), (5,2), (6,1)\}$.
\end{example}

\begin{example}[Infinite coin sequences]
$\Omega = \{H, T\}^{\mathbb{N}}$, an uncountable set. This is the natural
sample space for asymptotic statements such as the strong law of large
numbers.
\end{example}

\section*{3. The Power Set Theorem}

\begin{theorem}[Power Set Cardinality]
If $\Omega$ is a finite set with $|\Omega| = n$, then the power set
$2^{\Omega}$ has cardinality $2^n$.
\end{theorem}

\begin{proof}
We proceed by induction on $n$.

\medskip
\noindent\textbf{Base case} ($n = 0$). The unique set of cardinality $0$
is $\Omega = \emptyset$. Its only subset is $\emptyset$ itself, so
$|2^{\emptyset}| = 1 = 2^0$.

\medskip
\noindent\textbf{Inductive step.} Assume the claim holds for all sets of
cardinality $n$. Let $\Omega$ have cardinality $n + 1$. Fix any element
$x \in \Omega$ and let $\Omega' = \Omega \setminus \{x\}$, which has
cardinality $n$.

Partition the subsets of $\Omega$ into two classes:
\begin{enumerate}
\item Subsets that do \emph{not} contain $x$ --- these are exactly the
subsets of $\Omega'$, and by the inductive hypothesis there are $2^n$
of them.
\item Subsets that \emph{do} contain $x$ --- each such subset has the form
$A \cup \{x\}$ where $A \subseteq \Omega'$. The map $A \mapsto A \cup \{x\}$
is a bijection between $2^{\Omega'}$ and $\{B \subseteq \Omega : x \in B\}$,
so there are again $2^n$ such subsets.
\end{enumerate}
The two classes are disjoint and exhaust $2^{\Omega}$, so
\[
|2^{\Omega}| \;=\; 2^n + 2^n \;=\; 2 \cdot 2^n \;=\; 2^{n+1}.
\]
This completes the induction.
\end{proof}

\begin{corollary}
A coin flip ($n=2$) admits $2^2 = 4$ events; a die ($n=6$) admits
$2^6 = 64$ events; two dice ($n=36$) admit $2^{36} \approx 6.87 \times
10^{10}$ events.
\end{corollary}

\section*{4. Remarks on the Infinite Case}

When $\Omega$ is uncountable (e.g.\ $\Omega = [0,1]$), one cannot in
general use the full power set $2^{\Omega}$ as the event family --- there
exist subsets to which no translation-invariant probability measure can
be consistently assigned (the Vitali construction). Instead one restricts
to a $\sigma$-algebra $\mathcal{F} \subseteq 2^{\Omega}$, typically the
Borel $\sigma$-algebra. The triple $(\Omega, \mathcal{F}, \mathbb{P})$
is then called a \emph{probability space}; sample spaces are the first
slot of this triple. See node 18-sigma-algebras for the continuation.

\end{document}
